\begin{figure}[ht]
\centering
\begin{tikzpicture}[>=Stealth, font=\small, node distance=1.5cm]
  \node[draw, rounded corners, fill=blue!8, minimum width=3.4cm, minimum height=0.95cm] (count)
    {数点 $\#E(\F_p)$};
  \node[draw, rounded corners, fill=green!8, right=of count, minimum width=3.6cm, minimum height=0.95cm] (frob)
    {Frobenius 多项式};
  \node[draw, rounded corners, fill=orange!10, right=of frob, minimum width=3.6cm, minimum height=0.95cm] (L)
    {Euler 因子与函数方程};
  \node[draw, rounded corners, fill=purple!10, below=of frob, minimum width=4.2cm, minimum height=0.95cm] (bsd)
    {BSD：中心零点与秩};
  \draw[->, thick] (count) -- (frob);
  \draw[->, thick] (frob) -- (L);
  \draw[->, thick] (L) -- (bsd);
\end{tikzpicture}
\caption{本章习题的路线图：先从有限域点数读出 Frobenius 特征多项式，再进入局部 $L$-因子与函数方程，最后回到 BSD 所关心的秩与挠群。}
\label{fig:ch08-exercises-roadmap}
\end{figure}
